% !TeX root = ./main.tex
\section{Methodology}
\label{sec:methods}

\subsection{Problem formulation}
\label{sec:problem}

The forecasting task is posed as sequence-to-sequence regression on the time series of wave frequency spectrum measured at a single buoy. At each hourly time step $t$ the record contains five quantities on a fixed grid $f_1,\dots,f_F$ ($F=47$ bins, $0.0200$--$0.4850$~Hz, non-uniformly spaced with finer resolution at low frequencies): the spectral density $E(f,t)$ [m$^2$/Hz], the mean and principal wave directions $\alpha_1(f,t)$ and $\alpha_2(f,t)$, and the first and second normalised polar coordinates $r_1(f,t)$ and $r_2(f,t)$. The circular quantities $\alpha_1$ and $\alpha_2$ are encoded as $(\sin,\cos)$ pairs, giving $C=7$ channels per bin,
$\mathbf{x}(f_k,t) = [\,E,\ \sin\alpha_1,\ \cos\alpha_1,\ \sin\alpha_2,\ \cos\alpha_2,\ r_1,\ r_2\,](f_k,t)$.
 
Given the $L$ most recent hourly observations, $\mathbf{X}_t = \mathbf{x}(f_{1:F},\, t-L+1{:}t)$, the model forecasts a trajectory of $\tau$ hourly steps, $\tau \in \{12, 24, 48\}$~h, with a separate model trained and tuned for each horizon. Only the terminal step is treated as the forecast product; the intermediate steps are the autoregressive scaffolding needed to reach it.

\paragraph{Shape/magnitude decomposition of the forecast target}
Rather than forecasting the physical spectrum $E(f,t)$ directly, the forecast target is factorised into a shape and a magnitude component, which are modelled separately and recombined at inference:
\begin{equation}
E(f,t) \;=\; S(f,t)\, m_0(t), \qquad
S(f,t) = \frac{E(f,t)}{m_0(t)}, \qquad
m_0(t) = \Big(\tfrac{H_s(t)}{4}\Big)^{2},
\label{eq:shape_magnitude}
\end{equation}
where $S$ is the unit-area spectrum and $m_n = \int f^n E(f)\,\mathrm{d}f$ are the spectral moments. A shape model forecasts $S(f,\,t+1{:}t+\tau)$, a magnitude model forecasts $H_s$, and the two outputs are multiplied according to Eq.~\eqref{eq:shape_magnitude}. The split separates the comparatively easy magnitude problem from the harder shape problem, and is motivated by the tendency of a single model forecasting $E(f,t)$ directly to underestimate spectral peaks and over-smooth the high-frequency tail. Such a monolithic model is retained as a reference (Section~[INSERT: Results section reference]).
% MAGNITUDE BRANCH (kept as comments while the H_s model may stay out of the analysis):
% The two branches share the data pipeline, the encoder and the training schedule, but the
% magnitude model differs in four respects. Its decoder input is a scalar, embedded by a single
% linear layer rather than the frequency-structured embedding. Its loss is the RMSE of $H_s$ in
% metres, since the terms of Eq.~\eqref{eq:loss} are defined only for spectra. It is selected on
% the final-step skill score of $H_s$ against persistence rather than on PF. And
% $d_{\text{head}} \in \{8, 16, 32\}$ and $n_{\text{head}} \in \{4, 8\}$ are searched rather than fixed.

Spectral targets are predicted in log space, so the recovered spectrum is non-negative by exponentiation, without an output constraint or clipping. $T_{m02} = \sqrt{m_0/m_2}$, is derived from the forecast spectrum for evaluation while $H_s$ is output directly by the forecast model.

\subsection{Data}
\label{sec:data}

\paragraph{Source and split}
Data are from NDBC station 32012 (Woods Hole Oceanographic Institution Stratus mooring; $19^\circ25'29''$~S, $85^\circ04'39''$~W), which reports the five channels of Section~\ref{sec:problem} hourly. The record spans 2016-01-01 00:00 to 2017-12-31 23:00 (17,544 hourly observations, $\approx$2 years). It is split chronologically, without shuffling: 70\% for training (to 2017-05-26 15:00), 15\% for validation (to 2017-09-13 07:00) and 15\% for testing (to 2017-12-31 23:00). The test split plays no part in tuning or model selection, and final results are reported on it. Only temporal generalisation is assessed; generalisation to unseen buoys is not.
 
\paragraph{Preprocessing and normalisation}
All channels are reindexed to a common hourly grid, and gaps are filled by linear interpolation in time (bidirectional at the ends of the record). The directions are interpolated through their $\sin$/$\cos$ components and recovered with $\operatorname{atan2}$, so a gap is never bridged across the wrong side of the compass. Missing records are few and short. Of the 17,544 hourly records, 124 (0.71\%) were absent from the raw files and filled in this way. The five channels share identical timestamps, so each record is either fully observed or fully interpolated. The 124 missing records form 93 gaps, of which 87 last a single hour and only one exceeds 3~h (26~h, 30--31 May 2017, within the validation split). The model therefore sees a strictly regular sequence. Normalisation statistics are fitted on the training split only. The spectral density is scale-normalised, $\tilde{E}(f,t) = E(f,t)/\mu(f)$, with $\mu(f)$ the per-frequency training mean (floored at $10^{-8}$); scaling without centring preserves non-negativity, so $H_s = 4\sqrt{m_0}$ can be computed after denormalisation. $r_1$ and $r_2$ are $z$-scored, and the $\sin$/$\cos$ components are left unscaled.

\paragraph{Targets}
Targets are built from the physical density. The shape target $S = E/m_0$ (with $m_0$ floored at $10^{-12}$~m$^2$) is mapped to log space as $y = \log \max\!\big(y_{\text{phys}},\, \epsilon\,\bar{y}(f)\big)$, where $\bar{y}(f)$ is the per-frequency training mean of the physical target and $\epsilon = 10^{-3}$. The per-bin floor stays proportional to the energy each bin typically carries and avoids $\log 0$ in calm bins. The same transform is applied to the decoder start token and to every ground-truth decoder input, so the whole autoregressive loop runs in log space. Input/target pairs are formed by sliding a window of $L$ input and $\tau$ target steps across each split, giving $T_{\text{split}} - L - \tau + 1$ samples per split.
% MAGNITUDE BRANCH: the magnitude target is $H_s = 4\sqrt{m_0}$ of the physical density, in metres, with no normalisation or log transform.

\paragraph{Growth/decay features}
Whether the wave systems in the input window are growing or decaying is informative for their near-term evolution, but it is a property of the sequence of spectra and cannot be read from any single one. Each sample is therefore given twelve scalar features from Dynamic Mode Decomposition \citep{schmid2010dmd}. A linear operator $A$ with $\mathbf{x}_{t+1} \approx A\,\mathbf{x}_t$ is fitted across the $L$ normalised density spectra of the window, restricted to its eight leading singular vectors, and eigendecomposed. Each eigenvalue $\lambda$ gives a growth rate $\ln|\lambda|$ and an oscillation frequency $\arg\lambda$, both per hour. Because the spectra are real, eigenvalues come in conjugate pairs, and only the member with non-negative frequency is kept; the rank of eight, twice the number of modes retained, leaves enough candidates after this filtering. For the four remaining modes with the largest amplitude (the magnitude of their coefficient for the first spectrum of the window) we keep the growth rate, frequency and amplitude, with zeros where fewer than four modes exist. The features are $z$-scored on the training split and enter the encoder only.

\subsection{Model architecture}
\label{sec:architecture}
 
The model is an encoder--decoder Transformer \citep{vaswani2017attention} in the pre-normalisation arrangement \citep{xiong2020prenorm} (Fig.~\ref{fig:architecture}).
 
\begin{figure}[t]
  \centering
  % [INSERT: architecture schematic]
  \caption{Model architecture. Left: the four-stage frequency-structured embedding that maps the $F \times C$ measurement at one time step to a single token. Right: encoder--decoder layout, showing where the growth/decay features and the temporal convolutional front-end enter the encoder, and the autoregressive decoder loop with unit-area renormalisation.}
  \label{fig:architecture}
\end{figure}

\paragraph{Frequency-structured embedding}
Instead of flattening the $F \times C$ measurement at each time step through one linear layer, the embedding follows its (frequency, channel) structure in four stages:
\begin{enumerate}
  \item \emph{Per-bin projection.} A linear layer shared across bins maps the $C$ channels at each bin to an 8-dimensional representation (GELU). The channels at one frequency describe the same process, so they are combined before anything is mixed across frequencies.
  \item \emph{Frequency identity.} A sinusoidal encoding of each bin's log-frequency, scaled to the span of $\log f$, is added, together with a learned per-bin residual initialised at zero. Encoding the frequency value instead of the bin index matters because the grid is non-uniform: after a first gap of 0.0125~Hz, the spacing is 0.005~Hz up to 0.0925~Hz, 0.01~Hz from 0.1 to 0.35~Hz and 0.02~Hz from 0.365~Hz upwards.
  \item \emph{Frequency-axis convolution.} The dilated convolutional block described below is applied along frequency, with replicate padding, so each bin takes in context from its neighbours. Zero padding would introduce fictitious zero-energy bins beyond the grid edges. The block has the same design as the temporal one but its own weights.
  \item \emph{Attention pooling.} A learned query attends over the $F$ bins and pools their 8-dimensional representations into one summary vector, which a linear layer maps to the $d_{\text{model}}$-dimensional token. The weights depend on content, so they can follow the spectral peak, and the parameter count does not depend on $F$.
\end{enumerate}

\paragraph{Encoder}
The twelve growth/decay features are mapped jointly to $d_{\text{model}}$ by one linear layer, and the result is added to every encoder token. A fixed sinusoidal positional encoding \citep{vaswani2017attention} is then added; because the sequence is strictly regular, the step index suffices. Before self-attention, the encoder sequence passes through three dilated 1-D convolutions (kernel 3, dilations 1, 2, 4, length-preserving padding), each in a pre-norm residual block with GELU and dropout, so that local trends over a few hours are built into the tokens and attention can concentrate on longer-range dependencies.
% CHECK: receptive field with kernel 3, dilations 1,2,4 is 15 steps, not 9; symmetric padding makes the convolution centred, not causal.

\paragraph{Decoder}
Decoder inputs pass through the same four-stage embedding, instantiated with $C=1$, and receive the same positional encoding. The decoder is seeded with the last spectrum of the input window, converted to the target quantity and log-transformed like the targets. During inference each subsequent input is the model's previous prediction; during training a causal mask restricts attention to earlier positions. A final linear layer maps each position to a log-spectrum in $\mathbb{R}^F$. At inference, each predicted step is exponentiated, renormalised to unit area (trapezoidal rule on the physical grid) and returned to log space as in the targets before it is fed back, so the autoregressive state remains a valid density. In training, the prediction fed back under scheduled sampling (Section~\ref{sec:training}) is the raw log-space output, so the unit-area constraint is imposed only at inference. The encoder runs once per sample.

\paragraph{Capacity}
The model dimension is constructed as $d_{\text{model}} = d_{\text{head}} \times n_{\text{head}}$. For the shape model $d_{\text{head}}=32$ and $n_{\text{head}}=8$ are fixed ($d_{\text{model}}=256$), because they were consistently selected in earlier searches (Supplementary Section~S4). The freed search budget goes to the loss weights.

\subsection{Training}
\label{sec:training}
 
\paragraph{Loss function}
The spectral loss combines three terms evaluated on the log-space prediction $\hat{y}$ and target $y$,
\begin{equation}
\mathcal{L} \;=\; 
\lambda_{\text{KL}}\,D_{\text{KL}}
\;+\; \lambda_{W}\,W_2
\;+\; \lambda_{P}\,\mathcal{L}_{\text{peak}} .
\label{eq:loss}
\end{equation}

\begin{itemize}
  \item $D_{\text{KL}}(P_{\text{true}}\,\|\,P_{\text{pred}})$ is the Kullback--Leibler divergence \citep{kullback1951} between the spectra treated as distributions over frequency, $\log P = \operatorname{log\_softmax}(y + \log \Delta f)$. Because it is weighted by $P_{\text{true}}$, it charges mainly for mass missing where the true spectrum is concentrated, and a peak that broadens slightly is penalised far less than under a per-bin error.
  \item $W_2$ is the 1-D Wasserstein-2 distance between the mass-normalised spectra, computed in closed form in the quantile domain,
  $W_2^2 = \int_0^1 |F^{-1}(q)-G^{-1}(q)|^2\,\mathrm{d}q$,
  without an optimal-transport solver \citep{peyre2019}. It tolerates small shifts of a peak but charges a real transport cost for merging two peaks into one, which a pointwise loss does not distinguish from a large displacement.
  \item $\mathcal{L}_{\text{peak}}$ is a differentiable peak-height term. For each trough-to-trough window $k$ of the \emph{true} spectrum (Section~\ref{sec:evaluation}), the predicted peak height is a softmax-weighted average $\hat{H}_k = \sum_{i} \hat{E}(f_i)\,\sigma_k(i)$ with $\sigma_k(i) \propto \exp(\hat{E}(f_i)/\kappa_k)$, and the loss is the mean of $(\hat{H}_k - H_k^{\text{true}})^2$, where $H_k^{\text{true}}$ is the exact window maximum. The temperature $\kappa_k = \max\!\big(\kappa_{\min},\, c\,H_k^{\text{true}}\,(f_{r_k}-f_{l_k})/(f_F - f_1)\big)$, with $c=0.15$, scales with the peak's energy and bandwidth in Hz, so narrow swell peaks get a sharp softmax and broad wind-sea peaks a softer one. Because the windows are fixed by the true spectrum, the term ignores shifts of a peak within its own partition, and so complements $W_2$, which measures position but not amplitude. A peak displaced beyond its partition is penalised as missing, as it also is by $W_2$.
\end{itemize}
The weights are tuned per horizon. Implementation details are given in Supplementary Section~S2.

The loss ablation (Section~\ref{sec:ablation}) compares Eq.~\eqref{eq:loss} with a per-bin reference, the frequency-weighted mean squared error in log space, $\mathcal{L}_{\text{base}} = \sum_i w_i\,(\hat{y}_i - y_i)^2$, with trapezoidal quadrature weights $w_i \propto \Delta f_i$ summing to one. It enters with weight $\lambda_{\text{base}}$, which is zero in the final configuration (Table~\ref{tab:hyper}).

All terms take the log-space quantity the network emits as input, so no denormalisation step enters the gradient path. The per-bin loss weights each bin by its width rather than by the energy it carries, so a log-space error in a near-empty tail bin costs at least as much as the same error at the peak, although it barely affects $H_s$ or $T_{m02}$. $D_{\text{KL}}$ instead weights each bin's log-space error by the fraction of the true energy in that bin, and $\mathcal{L}_{\text{peak}}$ acts only within the detected peak windows.

\paragraph{Optimisation}
We use AdamW \citep{loshchilov2019adamw} with a linear learning-rate warmup that rises from 28\% to 100\% of the sampled rate over the first 5 epochs, so that the first epoch already makes progress, gradient clipping at norm 1.0, and learning-rate halving after 5 epochs without improvement (2-epoch cooldown, suspended during warmup). Without the warmup, the learning rate dominated hyperparameter importance in the search yet correlated only weakly with the objective, with the best and worst trials drawing it from much the same range; this points to unstable early training rather than an optimum the sampler could exploit (Supplementary Section~S1). Training uses scheduled sampling \citep{bengio2015scheduledsampling}: at each decoder step the ground-truth previous token is used with probability $p$, drawn per sample, and the model's detached prediction otherwise, with $p$ decaying linearly from 1 to 0 over the first 40 epochs. Training runs for at most 100 epochs, with early stopping after 20 epochs without improvement, and the best-validation checkpoint is kept. Because the validation objective is computed by autoregressive rollout on a short split and is noisy, every decision based on it (scheduling, checkpointing, early stopping and pruning) uses a 5-epoch trailing mean. All settings are listed in Table~\ref{tab:hyper}.

\subsubsection{Hyperparameter search}
\label{sec:hpo}

Hyperparameters (Table~\ref{tab:hyper}) are tuned with the multivariate Tree-structured Parzen Estimator \citep{bergstra2011tpe} in Optuna \citep{akiba2019optuna}, with 18 random startup trials and 80 trials per horizon. From epoch 30, and every 5 epochs thereafter, a median pruner compares a trial's best smoothed score with the median of completed trials at the same epoch, once at least five trials have completed. It is consulted only when the trial's own score has not improved on its earlier best for 10 epochs, since trials often dipped below the median around epoch 30, while teacher forcing was still being withdrawn, and then recovered (Supplementary Section~S1). The objective (Section~\ref{sec:evaluation}) is computed on the validation split.
The global seed is set once before the search and not reset per trial, so that trials differ in initialisation and data order and the sampler can separate hyperparameter effects from run-to-run noise. The selected configuration for each horizon is retrained from scratch with five seeds and evaluated once on the test split; we report mean $\pm$ standard deviation. Training used [INSERT: hardware/GPU specification].

\begin{table}[t]
\centering
\caption{Training settings and hyperparameter search space. Log-uniform ranges are marked (log).}
\label{tab:hyper}
\small
\begin{tabular}{ll}
\toprule
\multicolumn{2}{l}{\emph{Searched}} \\
Input window $L$ & $\{6,12, 24, 48, 96\}$~h \\
Batch size & $\{32, 64, 128\}$; $\{32, 64\}$ for $\tau > 24$~h \\
Learning rate & $10^{-3}$--$1.5\times10^{-2}$ (log) \\
Weight decay & $10^{-6}$--$10^{-2}$ (log) \\
Dropout, per-bin representation & $[0.1, 0.3]$ \\
Dropout, $d_{\text{model}}$-wide layers & $[0.0, 0.3]$ \\
Encoder / decoder layers & $\{1,2,3,4\}$ each \\
$\lambda_{\text{KL}}$, $\lambda_{W}$, $\lambda_{P}$ & $[0.1, 50]$, $[1, 200]$, $[0.01, 20]$ (log) \\
\midrule
\multicolumn{2}{l}{\emph{Fixed}} \\
$d_{\text{head}}$, $n_{\text{head}}$, $d_{\text{model}}$ & 32, 8, 256 \\
Per-bin width & 8 \\
Convolutions & kernel 3, dilations 1, 2, 4 \\
$\lambda_{\text{base}}$ & 0 \\
Warmup; gradient clip & 5 epochs (28--100\%); norm 1.0 \\
LR on plateau & factor 0.5, patience 5, cooldown 2 \\
Teacher forcing & $p$: $1 \to 0$ over 40 epochs, per-sample draw \\
Max epochs; early stopping & 100; patience 20 \\
Validation smoothing & 5-epoch trailing mean \\
Pruner (epochs) & median from 30, every 5; patience 10 \\
Log floor $\epsilon$ & $10^{-3}$ \\
Peak temperature $c$; $\kappa_{\min}$ & 0.15; $10^{-4}$ \\
Peak windows per spectrum & $\leq 4$ (lowest frequencies) \\
DMD modes & 4 \\
\bottomrule
\end{tabular}
\end{table}

\subsection{Evaluation}
\label{sec:evaluation}
 
\paragraph{Baselines}
Persistence repeats the last observed spectrum over all $\tau$ steps. A stronger baseline is a per-bin linear autoregressive model, rolled out recursively like the Transformer. Each bin is regressed on its own previous $p$ values and an intercept by least squares on the training split, with a ridge penalty of $10^{-6}$ that only stabilises the solve, since adjacent lags are highly correlated. It works in linear rather than log space on the same target as the Transformer, and for the shape target each predicted step is renormalised to unit area before it is fed back. The order is chosen independently of the Transformer's input window, from the same candidates $p \in \{12, 24, 48, 96\}$, on the validation split. The selection score is the per-step skill score against persistence, $1 - \mathrm{RMSE}/\mathrm{RMSE}_{\text{pers}}$ with trapezoid-weighted RMSE, averaged over the forecast steps with weights that halve every $\tau/2$ steps. It has no cross-frequency coupling and no nonlinearity, so it isolates the skill that comes from those two sources.
% CHECK: the AR order is selected on a skill score, not on PF like the Transformer, so the two are tuned on different criteria.

\paragraph{Core metrics and model selection}
We use a peak-fidelity criterion,

\begin{equation}
\text{PF} \;=\; F_1\!\left(\text{precision},\, \overline{\text{recall}}\right) \;\times\; \bigl(1 - \min(\overline{\text{rel.\,err.}},\, 1)\bigr),
\end{equation}
where $\overline{\text{recall}}$ is the class-averaged peak-separation recall, precision is the pooled peak-separation precision, and $\overline{\text{rel.\,err.}}$ is the class-averaged peak-height relative error (all defined below). Recall is averaged over the two classes so that a model cannot ignore whichever is rarer; precision is pooled, since a predicted peak that matches no true partition has no class to be assigned to. Both factors lie in $[0,1]$, so $\text{PF} \in [0,1]$ and neither can be traded away for the other. Recall alone is maximised by over-segmenting, which the precision term penalises; bounding the height error prevents a single badly-missed peak from dominating a criterion whose other term is a fraction. PF is not a skill score: no baseline enters it, so it is interpreted only by comparison against other forecasts scored identically.

\paragraph{Bulk and shape diagnostics}
In physical units and with trapezoidal moment integrals on the true grid, we report:
\begin{itemize}
  \item RMSE, bias and MAPE of $H_s$ for the recombined spectrum, pooled over all $\tau$ forecast steps (MAPE is well behaved here because $H_s \gtrsim 0.8$~m throughout the record);
  \item RMSE and bias of $T_{m02}$ on the final step, which is scale-invariant and can therefore be computed directly on the shape forecast;
  \item the per-bin scatter index $\mathrm{SI}(f) = \mathrm{RMSE}(f)/\overline{E_{\text{true}}(f)}$ for the recombined spectrum, pooled over all $\tau$ forecast steps, with the per-bin RMSE.
\end{itemize}
% CHECK: H_s and SI(f) of the recombined spectrum are pooled over all forecast steps, which conflicts with treating only the terminal step as the forecast product (Section problem). Restricting that evaluation to the final step would remove the conflict.

\paragraph{Peak-resolved, partition-conditioned diagnostics}
Because aggregate errors dilute errors confined to narrow peaks (Section~[INSERT: Introduction reference]), a set of peak-resolved metrics is computed on the final step. Peaks are detected with the significance criteria of \citet{portilla2009}: a local maximum is rejected if its frequency exceeds 0.4~Hz, its trough-to-trough partition holds less than 5\% of the total energy, it has fewer than two bins on either side before the next trough, or it lies between two higher-energy peaks. Each partition is labelled wind sea if $\gamma^{*} = E_{\text{obs}}(f_p)/E_{\text{PM}}(f_p) > 1$, where $E_{\text{PM}}$ is the Pierson--Moskowitz spectrum at the peak frequency; this peak-ratio criterion was proposed by \citet{portilla2009}, building on the spectral-fitting wind-sea identification of \citet{violantecarvalho2002}. Partitions and labels are taken from the true spectrum and applied to both spectra, so a partition the model misses counts directly as a recall miss. The metrics are:
\begin{itemize}
  \item the mean number of detected peaks, true and predicted; over-smoothing lowers the predicted count;
  \item peak-height relative error $|\hat{E}(f_p) - E(f_p)|/E(f_p)$ at each true peak, whether or not the model forms a peak there;
  \item peak-separation recall, the fraction of true peaks matched by a significant predicted peak within two bins;
  \item peak-separation precision, the fraction of predicted peaks that match a true peak within two bins. Recall counts only the peaks the model finds, so on its own it rewards a forecast that emits peaks indiscriminately; precision is what charges it for the ones that correspond to nothing. It is pooled rather than split by class, because a predicted peak with no true counterpart has no partition from which to inherit a label;
  \item $T_{m02}$ RMSE and bias within each true partition, excluding partitions whose predicted energy is below 5\% of the true energy.
\end{itemize}
All but precision and the peak counts are reported overall and separately for wind sea and swell. Wind-sea partitions are broad and fast-evolving, which makes them mainly a magnitude problem; swell partitions are narrow and persistent, which makes them mainly a position problem. A change that helps only one class can disappear in the pooled figure.
 
The normalisation round trip, the unit-area property of the shape transform and the stated properties of each loss term are checked with unit tests on synthetic spectra (Supplementary Section~S5; code in [INSERT: repository]).

\subsection{Incremental ablation strategy}
\label{sec:ablation}
 
Components were introduced one at a time and kept only if they improved the evaluation criteria in use at that stage without degrading others, at a computational cost justified by the gain. Because most candidates interact with capacity and regularisation, a candidate is not inserted into a frozen configuration: the full search of Section~\ref{sec:hpo} is re-run with it, under the same objective and search space as the comparison run, and the best configurations with and without it are compared. This mixes the candidate's effect with search noise, so no difference is claimed unless it exceeds the spread across five seeds of the selected configuration.
 
For the loss terms, where the candidate set is small but interaction with architecture is likely, the design was inverted. Seven arms (the per-bin baseline; each distributional term alone; KL combined with $W_2$ and with $\mathcal{L}_{\text{peak}}$; and all three) were run at the [INSERT: horizon] horizon with the architecture and training settings fixed, searching only the term weights; the baseline arm was repeated five times to estimate noise. Results are given in Section~[INSERT: Results section reference]. The computational cost of each component and the candidates that were not retained are reported in Supplementary Sections~S3 and~S4.
 
The protocol has two limitations. Attribution to a single component is only as reliable as the noise estimate, and differences of the order of the seed spread are reported as such. And the peak-resolved diagnostics were introduced partway through, so earlier components were accepted on RMSE-family metrics alone, which can favour over-smoothed forecasts, and have not been re-tested against the full panel.
